\section{Lesson 4: What Actually Follows?}

You can now state a passage's \textbf{Claim}, fix its \textbf{Meaning}, sort each premise's \textbf{Warrant}, and underline its \textbf{Burden}. Yet even after granting a burden-bearing premise, another trap waits: importing a tradition's destination into an early conclusion that has not earned it.

Bounding what an argument earns is the work of the fifth question---\textbf{Consequence}, paired with its conspicuous \textbf{Not yet established} note---under the rule that \emph{prose is canonical}: notation enters only when we name the exact feature it exposes.

\subsection{Validity, Soundness, and Bounding a Conclusion}

One rule governs inference: \emph{validity never upgrades warrant}. Validity preserves truth; it never elevates the kind of claim its premises make. Even a sound argument---valid with true premises---delivers only what its premises license. Because definition is not discovery, a valid deduction from definitional premises yields only a definitional conclusion; flawless form never turns verbal stipulation into discovery about reality.

\textbf{Consequence}---what follows if the burden is granted---has two halves: \emph{what follows}, stated at the weakest strength the premises support, and \emph{what is not yet established}, stated conspicuously beside it. Hold the first half to four bounds:
\begin{itemize}
\item \textbf{Scope:} all things of a kind, the cases at hand, or at least one? Arguments about changing things say nothing alone about unchanging things.
\item \textbf{Modality:} that something \emph{is} so, or \emph{must} be so? Premises about what happens establish what happens; necessity demands necessary premises.
\item \textbf{Subject matter:} reality, or our concepts?
\item \textbf{Uniqueness:} \emph{at least one}, or \emph{exactly one}? Moving from ``a first'' to ``the first'' normally needs a separate argument.
\end{itemize}

Test any conclusion by weakening each clause in turn---trimming scope, lowering ``must'' to ``is,'' dropping ``the one'' to ``at least one''---to see what survives. A bounded result may stand as established, conditional on an open burden, conceptual, or not yet established---a completed finding, never a failure.

The second half, \textbf{Not yet established}, holds open the gap to the usual destination. ``There is a first mover in this order of motion'' is not ``there is one God, pure act, creating from nothing.'' Generic cautions (``much remains to be said'') are empty; name the specific unproved claims---uniqueness, immateriality, simplicity, intelligence, goodness, creation---so later arguments can be read as arguments.

\subsection{Seven Recurring Overreaches and the Discipline of Scope}

Conclusions strengthen quietly along seven recurring paths. Watch for them in ordinary English as checkpoints rather than automatic verdicts, since explicit premises sometimes license the step: \emph{from is to must} (actuality upgraded to necessity); \emph{from these to all} (cases at hand stretched universally); \emph{from a concept to a thing} (a notion's requirement treated as existing); \emph{from at least one to exactly one} (``a cause'' turned into ``the sole cause''); \emph{from that to why} (showing \emph{that} something is so read as explaining \emph{why}, Lesson 6's \emph{quia} and \emph{propter quid} boundary); \emph{from a cause to this cause} (identifying some cause with a familiar candidate without proof); and \emph{from the parts to the whole} (transferring a member's trait directly to the collection).

Before any notation appears, catch three scope traps in ordinary English.

First, \textbf{quantifier scope}: ``everything moved has some mover'' is not ``there is something that moves everything moved.'' The first gives each moved thing its own mover; the second sets one mover over all moved things. Just as ``everyone has a mother'' never entails ``one woman is everyone's mother,'' the first never entails the second.

Second, \textbf{modal scope}: ``necessarily, if this is moved, it has a mover'' is not ``if this is moved, it has a necessary mover.'' The first attaches necessity to the whole conditional (\emph{de dicto}, about the statement); the second attaches it to the mover itself (\emph{de re}, about the thing). That slide drives fatalism and careless readings of proofs for a necessary being.

Third, \textbf{composition and division}: moving from part to whole, whole to part, or every member to the collection. Every brick is clay, so the wall is clay; yet every brick weighs under four pounds while the wall does not, and an ancient guild need not have ancient members. Whenever an argument crosses between members and wholes, ask what licenses the crossing.

\subsection{Prose Is Canonical: When Notation Helps and When It Hides}

Prose is the canonical medium of metaphysical reasoning. Notation never appears by default, never projects pseudo-rigour, and is never required of the reader. It earns its place only for seven features: \textbf{quantifier scope}, \textbf{modal scope}, \textbf{contradiction}, \textbf{identity}, \textbf{simple dependency}, a \textbf{formally invalid inference}, or \textbf{two genuinely different formal readings} of one sentence. Keep one rule: \emph{symbolize the logic around the metaphysics, not the metaphysics itself, unless the latter genuinely becomes clearer.}

Every displayed formula requires an explicit \textbf{licensing sentence} right beside it naming the exact feature it exposes; without that sentence, cut it. Hence \textbf{act, potency, participation, essence, existence, analogy, form, matter, causation, and grounding} are \emph{not} routinely encoded as predicate letters. Their live disputes ask what each term means and whether it names anything real---questions predicate letters conceal while making them look settled. Routine symbolization produces two failures.

\textbf{Failure A: symbols that hide the question.} This formula is displayed strictly as a counterexample exposing how predicate letters conceal open metaphysical questions inside tidy symbols (\linline{\lall} means ``for every''; \linline{\limp} means ``if \dots\ then''):
\ldisplay{\lall x\lthin (\logicop{Change}(x) \limp \logicop{Potency}(x))}
Well-formed yet useless, it packs every live dispute from Lesson 3 inside two letters: what \emph{change} and \emph{potency} mean, whether potency is a real principle or a description of future states, and why change licenses that claim. Whenever the open question sits inside a predicate letter, drop the symbols and argue in prose.

\textbf{Failure B: tautology dressed as discovery.} Formalizing definitions that already guarantee a conclusion lends polish where substance is absent. Keep \emph{definitional consequence} apart from \emph{substantive metaphysical conclusion} via the hostile-reader test: if no opponent can accept the definitions and deny the conclusion, symbols have dressed vocabulary as discovery. Often, as in Worked Passage 2, the clearest analysis uses zero symbols.

\subsection{Worked Passage 1: Bounding the First Way}

We return to St.\ Thomas Aquinas's First Way in \emph{Summa Theologiae} I, q.\ 2, a.\ 3, \emph{corpus} (1920 English Dominican translation, checked against Leonine Latin, Rome, 1888; parallel at \emph{Summa contra Gentiles} I, c.\ 13, n.\ 14). Lesson 3 completed \textbf{Claim}, \textbf{Meaning}, \textbf{Warrant}, and \textbf{Burden}; here we bound \textbf{Consequence} alone, leaving \textbf{Denial} for Lesson 5:

\begin{quote}
Now it is not possible that the same thing should be at once in actuality and potentiality in the same respect, but only in different respects. [\dots] Therefore, whatever is in motion must be put in motion by another. [\dots] Therefore it is necessary to arrive at a first mover, put in motion by no other; and this everyone understands to be God.
\end{quote}

\subsubsection*{Consequence}

Grant Lesson 3's twofold \textbf{Burden}: real potency requiring a distinct actual mover, and subordinate chains like staff and hand. What actually follows?

First, two formulas earn their place by exposing a quantifier-order gap: every moved thing having some mover other than itself versus one thing moving everything moved (\linline{\lex} means ``there is some''; \linline{\landop} means ``and''; \linline{M(x)} abbreviates ``\linline{x} is moved'' and \linline{\logicop{Moves}(y, x)} ``\linline{y} moves \linline{x},'' setting aside the ontology of motion strictly to show scope):
\ldisplay{\lall x\lthin (M(x) \limp \lex y\lthin (y \lneqop x \landop \logicop{Moves}(y, x)))}
\ldisplay{\lex y\lthin \lall x\lthin (M(x) \limp \logicop{Moves}(y, x))}
The premise supplies the first line; a single shared cosmic mover needs at least the second, and strict uniqueness requires a further argument still. Even excluding endless subordinate chains yields only at least one unmoved mover per chain---not one common mover over all motion, still less a unique one. Jumping from the first line to a unique cosmic mover commits the leap \emph{from at least one to exactly one}.

Second, two modal displays earn their place by exposing two modal-scope limits: Aquinas bars only simultaneous act and potency in one respect (not contingency in a changing thing), and necessity over a conditional never yields a necessary mover (\linline{\lneg} means ``not''; \linline{\ldiam} ``possibly''; \linline{\lbox} ``necessarily''; \linline{A(x)} and \linline{P(x)} denote simultaneous act and potency in the same respect, setting aside the ontology of potency to show scope):
\ldisplay{\lneg \ldiam (A(x) \landop P(x))}
Impossibility governs solely that simultaneous conjunction, never claiming a changing thing could not have been otherwise.

\keeppair
Compare next where necessity sits in ``whatever is in motion must be put in motion by another'':
\ldisplay{\lbox (M(x) \limp \lex y\lthin \logicop{Moves}(y, x))}
\ldisplay{M(x) \limp \lex y\lthin \lbox \lthin \logicop{Moves}(y, x)}
The upper line places necessity over the conditional; the lower pushes it onto the mover itself. Aquinas's premise licenses at most the upper line, never a mover that moves or exists necessarily.

Weakening each clause leaves the bounded conclusion: \emph{granting both burdens, there is at least one mover, not itself moved in the respect at issue, at the head of that essentially ordered series.}

\textbf{Not yet established.} This paragraph does not establish that there is only \emph{one} such mover, nor that it is immaterial, simple, pure act in every respect, intelligent, good, or Creator from nothing, nor that it is God. ``And this everyone understands to be God'' is a naming; Aquinas argues entitative attributes separately at \emph{ST} I, qq.\ 3--11, divine knowledge at q.\ 14, and across \emph{SCG} I, cc.\ 14--28.

\subsection{Worked Passage 2: A Dependence Claim Without Symbols}

Our second passage is St.\ Thomas Aquinas, \emph{De ente et essentia} (\emph{On Being and Essence}), Chapter 4 (Corpus Thomisticum Latin text, based on Baur and Roland-Gosselin):

\begin{quote}
Omnis autem essentia vel quiditas potest intelligi sine hoc quod aliquid intelligatur de esse suo; possum enim intelligere quid est homo vel Phoenix et tamen ignorare an esse habeat in rerum natura. Ergo patet quod esse est aliud ab essentia vel quiditate, nisi forte sit aliqua res, cuius quiditas sit ipsum suum esse [\dots] Ergo oportet quod omnis talis res, cuius esse est aliud quam natura sua habeat esse ab alio. Et quia omne, quod est per aliud, reducitur ad illud quod est per se sicut ad causam primam, oportet quod sit aliqua res, quae sit causa essendi omnibus rebus, eo quod ipsa est esse tantum.
\end{quote}

A close English rendering of the Latin:

\begin{quote}
Every essence or quiddity can be understood without anything being understood about its existence; for I can understand what a man or a phoenix is and yet not know whether it has existence in the nature of things. Therefore it is clear that existence is other than essence or quiddity, unless perhaps there is some reality whose quiddity is its very existence [\dots] Therefore it must be that every such thing, whose existence is other than its nature, has its existence from another. And because everything that is through another is reduced to that which is through itself as to a first cause, there must be some reality that is the cause of being for all things, because it is existence alone.
\end{quote}

\textbf{Zero symbolic notation} appears here---our clearest case of notation refused for a reason. Turning existence (\emph{esse}) into a predicate letter hides the very metaphysical question at stake: whether existence is the kind of thing a predicate letter can represent, or a real act composing with essence. Prose alone keeps that question open.

\subsubsection*{Claim}

Aquinas asserts that except in a reality whose essence is its existence, existence differs from essence; therefore every such reality receives existence from another, leading to a first cause that is pure existence and causes being in all things.

\emph{Plain-language reading.} Because we can grasp what a man or phoenix is without knowing whether any exists, \emph{what a thing is} differs from \emph{that it exists} in finite things; hence every finite thing receives existence from outside, ending in a first cause that simply is existence. We retain \emph{essence} and \emph{existence} so the shift from concept to real composition stays inspectable.

\subsubsection*{Meaning}

\emph{Essence} or \emph{quiddity} names what a thing is; \emph{existence} (\emph{esse}) names real being (\emph{in rerum natura}), not grammatical predication. \emph{Other than} (\emph{aliud ab}) means real composition inside the thing, not two mental concepts. \emph{Per aliud} (through another) contrasts received being with self-possessed being (\emph{per se}).

\subsubsection*{Warrant}

Conceiving a man or phoenix without knowing whether it exists rests on \textbf{experience} of our concepts. Stepping from conceivability to real distinction reads as a \textbf{logical consequence} (``Therefore it is clear''), yet leans silently on a \textbf{metaphysical principle} bridging mental separation to real composition. In the omitted lines between ``therefore'' clauses, Aquinas invokes a \textbf{metaphysical principle} (features flow from nature or outside) and a \textbf{logical consequence} barring efficient self-causation. Reducing dependent beings to pure existence invokes a \textbf{metaphysical principle} excluding an infinite regress of causes of being.

\subsubsection*{Burden}

Here the text stops describing a distinction and starts asserting a cause: grant real composition of essence and existence, and causal dependence follows quickly. The primary burden sits right at the start: \emph{that conceiving an essence without knowing whether it exists establishes a real distinction between essence and existence inside existing things}---assumed or crossed in one ``therefore'' here rather than argued against a conceptual distinction. A secondary burden is the unargued ban on an infinite regress of causes of being.

\subsubsection*{Consequence}

Granting the real distinction, the ban on self-causation, and the no-regress principle establishes that every reality whose existence differs from its nature depends efficiently for its being on at least one cause not receiving existence from another---whose essence is its existence.

\textbf{Not yet established.} Without auditing further steps in Chapter 4, this excerpt does not establish that this first cause is \emph{one}, simple in every further sense, intelligent, free, good, or provident, or creates freely rather than necessarily. Lesson 5 takes up \textbf{Denial}, where rejecting the real distinction proves a serious rival framework.

\subsection{Worked Passage 3: Member to Whole and the Burden Relocated}

A bad defence can be worse than the argument it defends, and errors belong to whoever made them. Consider this argument, \textbf{explicitly labelled as a reconstruction rather than Aquinas's words}:

\begin{quote}
\emph{Reconstruction (not Aquinas's text):} Every member of the series of moved movers is moved by another; therefore, the series as a whole is moved by another.
\end{quote}

Read it against St.\ Thomas Aquinas, \emph{Summa Theologiae} I, q.\ 46, a.\ 2, ad 7 (``Whether It Is an Article of Faith That the World Began?'', 1920 English Dominican translation, checked against Leonine Latin, Rome, 1888). In the body, Aquinas holds the world's temporal beginning strictly by faith; Objection 7 urged that an eternal world entails an impossible infinite regress of fathers begetting sons. Reply 7 answers:

\begin{quote}
In efficient causes it is impossible to proceed to infinity \emph{per se}---thus, there cannot be an infinite number of causes that are \emph{per se} required for a certain effect; for instance, that a stone be moved by a stick, the stick by the hand, and so on to infinity. But it is not impossible to proceed to infinity \emph{accidentally} as regards efficient causes; for instance, if all the causes thus infinitely multiplied should have the order of only one cause, their multiplication being accidental, as an artificer acts by means of many hammers accidentally, because one after the other may be broken. It is accidental, therefore, that one particular hammer acts after the action of another; and likewise it is accidental to this particular man as generator to be generated by another man; for he generates as a man, and not as the son of another man. For all men generating hold one grade in efficient causes---viz.\ the grade of a particular generator. Hence it is not impossible for a man to be generated by man to infinity [\dots]
\end{quote}

\subsubsection*{Claim}

The reconstruction claims that because every moved mover is moved by another, the series as a whole is moved by another. Aquinas claims far less: infinite regress is impossible in causes ordered \emph{per se}, yet not judged impossible (\emph{non reputatur impossibile}) in causes ordered \emph{per accidens}.

\emph{Plain-language reading.} The reconstruction transfers a trait of every link straight to the chain. Aquinas refuses that shortcut: subordinate tools acting on borrowed power cannot regress endlessly, whereas independent causes acting sequentially---broken hammers, or fathers begetting sons---could stretch backward without demonstrated impossibility.

\subsubsection*{Meaning}

\emph{Per se} causes are concurrently required for an effect, lower members acting only through higher ones (hand, stick, stone). \emph{Per accidens} causes share one causal grade, each acting by its own nature rather than its predecessor's present action. \emph{Member} names an item in a series; \emph{whole} names the collection as one subject.

\subsubsection*{Warrant}

The reconstruction pairs a \textbf{metaphysical principle} with an attempted \textbf{logical consequence}. In Reply 7, stick, hammers, and generations appeal to \textbf{experience}; the world's beginning stays labelled as \textbf{revelation}; and the ban on endless \emph{per se} causes enters as a \textbf{metaphysical principle} carried as a \textbf{prior result} argued at \emph{SCG} I, c.\ 13, nn.\ 14, 33---standing on \textbf{provenance only} here until checked.

\subsubsection*{Burden}

The reconstruction hides a false bridge from members to wholes. Reading Aquinas removes that fallacy yet \emph{relocates rather than removes the burden}: the weight now sits on whether \emph{real natural dependencies are ordered per se rather than per accidens, and whether an endless per se series produces no effect.}

\subsubsection*{Consequence}

One display earns its place here by exposing a formally invalid inference---from every member of collection \linline{S} having feature \linline{F} to \linline{S} itself having \linline{F} (\linline{x \lin S} reads ``\linline{x} is a member of \linline{S}''):
\ldisplay{\lall x\lthin (x \lin S \limp F(x))}
\ldisplay{F(S)}
The upper line does not entail the lower. One counterexample settles it: take three two-pound bricks; each weighs under four pounds, yet the three together weigh six pounds---not under four.

The crude reconstruction fails \emph{from the parts to the whole}, and \textbf{the error belongs to the reconstructor, not to Aquinas}. Granting Reply 7 yields only a conditional result: \emph{infinite regress is impossible in a per se series of efficient causes, yet not excluded in a per accidens series.}

\textbf{Not yet established.} Reply 7 does not establish that any cosmic chain beyond simultaneous tools satisfies \emph{per se} ordering, nor that separate \emph{per se} chains share one first cause, nor (as Aquinas insists) that the universe began in time by natural reason. The distinction saves the argument from a blunder while placing its real burden in clear view.

\subsection{Exercises}

\subsubsection*{Exercise 4.1}

Read this annotated argument and the proposed bound beneath it:

\begin{quote}
(\emph{Experience}) Every observed living organism begins to exist and perishes. (\emph{Metaphysical principle}) Whatever begins to exist requires a distinct originating cause. (\emph{Metaphysical principle}) An essentially ordered chain concurrently sustaining a thing's present existence cannot stretch to infinity. (\emph{Final sentence}) Therefore, there must exist one eternal, all-powerful Creator who made the universe from nothing.
\end{quote}

\emph{Proposed bound:} ``Granting these premises establishes one necessary, non-perishing cause producing all living organisms.''

Reject this proposed bound as \textbf{still too strong}, naming two unlicensed leaps inside it. Weaken the conclusion to what actually follows, and note under \textbf{Not yet established} the gap to the final sentence.

\subsubsection*{Exercise 4.2}

Read this short passage:

\begin{quote}
Every natural event is necessarily brought about by some prior physical state; without an antecedent cause, no event could occur. Since every event therefore depends on a necessary cause, nature's whole history unfolds by strict necessity, excluding contingency.
\end{quote}

Locate the hidden modal or quantifier shift and explain in ordinary English what it conceals.

\subsubsection*{Exercise 4.3}

Read this unannotated paragraph from a fluent critic of Aquinas's First Way:

\begin{quote}
Aquinas's First Way fails on a first-week logic blunder. From ``everything moved is moved by another,'' Aquinas leaps straight to one First Mover behind everything---like arguing that because everyone has a mother, one woman is everyone's mother. Moreover, an endless past chain of movers, like parents begetting children, involves no contradiction. Nothing in the argument survives.
\end{quote}

Explain where this critic misfires against Aquinas's texts (\emph{ST} I, q.\ 2, a.\ 3, \emph{corpus}, and \emph{ST} I, q.\ 46, a.\ 2, ad 7), and name the genuine uniqueness gap that remains.

\subsubsection*{Exercise 4.4}

Read this unannotated passage from St.\ Thomas Aquinas, \emph{Summa Theologiae} I, q.\ 2, a.\ 3, \emph{corpus} (the Third Way, 1920 English Dominican translation):

\begin{quote}
The third way is taken from possibility and necessity, and runs thus. We find in nature things that are possible to be and not to be, since they are found to be generated, and to corrupt, and consequently, they are possible to be and not to be. But it is impossible for these always to exist, for that which is possible not to be at some time is not. Therefore, if everything is possible not to be, then at one time there could have been nothing in existence. Now if this were true, even now there would be nothing in existence, because that which does not exist only begins to exist by something already existing. Therefore, if at one time nothing was in existence, it would have been impossible for anything to have begun to exist; and thus even now nothing would be in existence---which is absurd. Therefore, not all beings are merely possible, but there must exist something the existence of which is necessary. But every necessary thing either has its necessity caused by another, or not. Now it is impossible to go on to infinity in necessary things which have their necessity caused by another, as has been already proved in regard to efficient causes. Therefore we cannot but postulate the existence of some being having of itself its own necessity, and not receiving it from another, but rather causing in others their necessity. This all men speak of as God.
\end{quote}

Flag any scope shift, state the bounded conclusion if its premises are granted, and identify what is \textbf{not yet established}.

\subsubsection*{Exercise 4.5}

Read this unannotated passage from St.\ Thomas Aquinas, \emph{Summa Theologiae} I, q.\ 2, a.\ 3, \emph{corpus} (the Fifth Way, 1920 English Dominican translation). Nothing is supplied. Run \textbf{Claim} (with a plain-language reading), \textbf{Meaning}, \textbf{Warrant}, \textbf{Burden}, and \textbf{Consequence} (including \textbf{Not yet established}) unaided:

\begin{quote}
The fifth way is taken from the governance of the world. We see that things which lack intelligence, such as natural bodies, act for an end, and this is evident from their acting always, or nearly always, in the same way, so as to obtain the best result. Hence it is plain that not fortuitously, but designedly, do they achieve their end. Now whatever lacks intelligence cannot move towards an end, unless it be directed by some being endowed with knowledge and intelligence; as the arrow is shot to its mark by the archer. Therefore some intelligent being exists by whom all natural things are directed to their end; and this being we call God.
\end{quote}

\subsection{Solutions}

\subsubsection*{Solution 4.1}

The proposed bound stays too strong on three counts: it leaps \emph{from at least one to exactly one} (``one \dots\ cause''), leaps \emph{from is to must} (``necessary''), and slides \emph{from these to all} while merging originating and sustaining causes.

Weakening each clause leaves a conclusion markedly weaker than the final sentence: \emph{every observed living organism had at least one distinct originating cause, and any essentially ordered chain concurrently sustaining its present existence ends in at least one cause not concurrently sustained in that respect.}

\textbf{Not yet established.} Nothing here proves that there is only \emph{one} such cause, eternal, non-perishing, necessary, or all-powerful, extending to the universe, or creating from nothing.

\subsubsection*{Solution 4.2}

The hidden move is a \textbf{modal scope shift} in Sentence 2---an overreach \emph{from is to must}. Sentence 1 places necessity over the conditional link: necessarily, if an event occurs, some prior state caused it. Sentence 2 moves ``necessary'' inside the clause, rewriting ``necessarily has a cause'' as ``depends on a necessary cause.'' Contingent causes never make outcomes necessary. A secondary slide \emph{from the parts to the whole} enters in jumping to ``nature's whole history''; both notes hold, with the modal shift doing the main work.

\subsubsection*{Solution 4.3}

The critic misfires twice. First, in \emph{ST} I, q.\ 2, a.\ 3, \emph{corpus}, Aquinas inserts an ordered regress argument (staff and hand) rather than leaping straight from ``moved by another'' to a first mover. Second, an endless chain of parents begetting children attacks a claim Aquinas expressly \emph{rejects} at \emph{ST} I, q.\ 46, a.\ 2, ad 7, where infinite \emph{per accidens} succession is not judged impossible; he bars regress only in \emph{per se} causes.

Yet the \textbf{genuine uniqueness gap} survives: granting the \emph{per se} argument yields at least one unmoved mover per chain, not yet one First Mover over all motion. Aquinas argues divine unity separately across \emph{ST} I, qq.\ 3--11.

\subsubsection*{Solution 4.4}

Two scope shifts appear. First, moving from ``that which is possible not to be at some time is not'' to ``at one time there could have been nothing in existence'' shifts quantifier scope---from \emph{each contingent thing fails at some time} to \emph{one time when nothing exists}---without ruling out overlapping contingent chains. Second, the final sentence shifts from ``not receiving [necessity] from another'' to ``causing in others their necessity,'' adding universal causation beyond ending a regress.

Granting the substantive premises yields only: \emph{there exists at least one being whose necessity is not caused by another.}

\textbf{Not yet established.} The Third Way does not prove there is only \emph{one} self-necessary being (note ``some being''), nor that it causes necessity in all others, nor that it is immaterial, simple, intelligent, free, good, or Creator from nothing. ``This all men speak of as God'' is a naming; attributes follow at \emph{ST} I, qq.\ 3--11 and q.\ 14.

\subsubsection*{Solution 4.5}

\textbf{Claim.} Aquinas claims that because non-intelligent bodies consistently act in the same way for the best result, they act for an end designedly; since nothing lacking intelligence tends toward an end unless directed by an intelligent being (like arrow and archer), an intelligent being directs all natural things to their end, called God.

\emph{Plain-language reading.} Unthinking bodies reliably reach beneficial outcomes, showing goal-directedness rather than luck. Since an unthinking thing aims at a goal only if guided by a mind, an intelligent being directs all natural things, called God.

\textbf{Meaning.} \emph{Act for an end} means intrinsic orientation to a determinate outcome, not conscious intent. \emph{Designedly} contrasts directed regularity with chance; \emph{directed by} means guidance by an intelligent agent.

\textbf{Warrant.} Regular outcomes rest on \textbf{experience}; inferring intrinsic goal-directedness is a \textbf{metaphysical principle}; requiring an intelligent director is a second \textbf{metaphysical principle}, illustrated from \textbf{experience} by arrow and archer; the conclusion enters as a \textbf{logical consequence} plus a naming.

\textbf{Burden.} Two independent burdens carry the argument: \emph{that regular beneficial outcomes prove intrinsic goal-directedness rather than blind law}, and \emph{that non-intelligent goal-directedness requires an external intelligent director}. Naming either as primary is defensible, since denying either blocks the proof; naming both locates the whole burden.

\textbf{Consequence.} Granting both burdens yields only that \emph{every non-intelligent natural thing acting for an end is directed by at least one intelligent being}---not yet that one single mind directs \emph{all} natural things.

\textbf{Not yet established.} The Fifth Way does not establish that there is only \emph{one} directing intelligence, nor that it is infinite, uncreated, simple, all-powerful, morally good, or Creator of matter rather than an external designer. Having learned to bound what follows when a burden is granted, we turn in Lesson 5 to our final question: what happens when a reader says no?
